\section{Optimality}
\label{sec:optimality}

\gp{A troublesome optimality result}

\gp{Uniqueness + implications}

\gp{Can't do better than a random coin toss.}

In this section, we prove that the post-processing method \autoref{alg:method}
produces a solution which is not only feasible but also near-optimal over
\emph{all} non-discriminatory classifiers with respect to generalized
false-positive and false-negative rates. Intuitively, this holds because,
under the calibration constraint, the cost of a classifier uniquely determines
the false-positive/false-negative rates for a given base rate $\mu_t$:
%
\begin{lemma}
  Let $h_t$ and $h'_t$ be perfectly calibrated classifiers with cost
  $\fairconst{h_t}{t} \leq \fairconst{h'_t}{t}$ for some
  cost function $\fairconstname$. Then $\fp{h_t} \leq \fp{h'_t}$
  and $\fn{h_t} \leq \fn{h'_t}$, with equality only if $\fairconst{h_t}{t} = \fairconst{h'_t}{t}$.
  \label{lem:fpfn_exact}
\end{lemma}
%
This also holds in an approximate sense if $h_t$ and $h'_t$ are approximately
calibrated (see \autoref{lem:fpfn} in \autoref{sec:cost_and_error}).
As described in \autoref{sec:setup}, the false-positive rate of a calibrated classifier
determines the false-negative rate, and vice versa. Because cost functions
of the form given by \eqref{eqn:cost_general} are positive combinations of
the false-positive/false-negative rates, it can be seen that there is only one
false-positive/false-negative solution for a given cost value, and improving one
of the values strictly improves the others.
The geometric intuition of this can be seen in \autoref{fig:postprocess}:
the level-order curves of $\fairconstname$ intersect the set of
calibrated classifiers in one point on the false-positive/false-negative plane.

Our goal is to find classifiers that simultaneously minimize the
false-positive/false-negative costs for both groups $G_1$ and $G_2$, subject to a
calibration constraint and non-disparity.
We continue with the assumption that, given a finite set of data and features,
we cannot strictly improve the cost (and therefore the
false-positive/false-negative rates) of the discriminatory classifiers $h_1$ and $h_2$.
Therefore, non-discriminatory classifiers with cost lower than that of
$\fairconst{h_1}{1}$ are infeasible (again assuming without loss of
generality that $\fairconst{h_1}{1} \geq \fairconst{h_2}{2}$).
Note that \autoref{alg:method} produces a classifier $\fair h_2$ with cost equal
to $\fairconst{h_1}{1}$, therefore achieving the minimum possible
cost for non-discriminatory classifiers. If we assume that $h_1$ and $\fair h_2$
are perfectly calibrated, then by \autoref{lem:fpfn_exact} $h_1$ and $\fair h_2$
must also achieve the minimum possible false-positive and false-negative rates
for non-discriminatory classifiers, and thus are optimal. This is summarized below:
%
\begin{thm}[Optimality of \autoref{alg:method}]
  Let $h_1$ and $h_2$ be (discriminatory) classifiers which achieve the
  minimum false-positive and false-negative rates for groups $G_1$ and $G_2$, respectively
  (i.e. cannot be strictly improved).
  If $h_1$ and $h_2$ are perfectly calibrated, then \autoref{alg:method} produces
  the optimal perfectly calibrated non-discriminatory classifiers with respect
  to false-positive/false-negative rates.
  \label{thm:optimality}
\end{thm}
%
\autoref{thm:optimality} holds in an approximate sense as well:
if $h_1$ and $h_2$ are approximately (but not perfectly) calibrated, then
\autoref{alg:method} produces non-discriminatory classifiers which are
near-optimal among classifiers that are just as calibrated as $h_1$ and $h_2$.
We formalize this in \autoref{thm:approx-fp-optimality} and
\autoref{thm:approx-fn-optimality} in \autoref{sec:optimality-proofs}.

%\paragraph{Bounding Generalized False-Positive and False-Negative Rates.}
%We can show that our postprocessing technique is near-optimal with respect
%to generalized false-positive and false-negative rates.
%We assume that the given classifier $h_1$ is Pareto-dominant in the sense that
%we cannot find a strictly better classifier $h_1'$ that is at least as
%calibrated as $h_1$ but has strictly smaller generalized false-positive or false-negative
%rate. As a consequence, $\fairconst{h_1}{1}$ and $\fairconst{\fair
%h_2}{2}$ must both be at least $\fairconst{h_1}{1}$. (We formalize this with \autoref{lem:fpfn}
%in the supplementary material).
%As a result, assuming that $h_1$ was
%Pareto-dominant, any classifiers satisfying our calibration and disparity
%constraints must have generalized false-positive and false-negative rates
%approximately equal to those of our derived classifiers $h_1$ and $\fair
%h_2$.

\paragraph{Relation to Equalized Odds post-processing.}
It is worth noting that \autoref{thm:optimality} implies that
the classifiers produced by \autoref{alg:method} are optimal not
only for post-processing methods but among \emph{all}
calibrated classifiers.
\gp{How to word this better?}
%
This is in contrast this with the (uncalibrated) Equalized Odds
framework, in which post-processing methods do not always produce
the optimal non-discriminatory classifier.
\citet{woodworth2017learning} show that the post-processing method
of \citet{hardt2016equality} may be suboptimal when group membership
is more predictive than any other feature.
It is worth noting that this result does not directly apply to
\methodname{} non-discrimination.
Firstly, \methodname{} is infeasible when group membership is the
most predictive feature (\autoref{crly:infeasible_predictive}), and
therefore the suboptimality result of \cite{woodworth2017learning} does
not apply.
Moreover, the feasibility region of classifiers for \methodname{}
is sufficiently more restrictive than that of Equalized Odds, which
makes optimization through post-processing tractable.
Equalized Odds solutions may lie anywhere in the false-positive/false-negative
plane, as the framework affords two degrees of freedom. As the post-processing
method of \cite{hardt2016equality} confines solutions to a limited subregion of the
plane it is possible that the derived classifiers may not be optimal. However,
solutions to \methodname{} are restricted to a line in the false-positive/false-negative
plane. As \autoref{alg:method} can derive all possible classifiers in this subspace
(again assuming we cannot strictly improve upon $h_1$ and $h_2$), then a solution
solution which is optimal with respect to the post-processing method \autoref{alg:method}
must be optimal over all (achievable) non-discriminatory solutions.

\paragraph{Implications.}
The optimality of \autoref{alg:method} is beneficial for many reasons.
Beyond ease of implementation, post-processing algorithms offer visibility into
the performance/disparity trade-off, whereas training-time constraints make
this tradeoff harder to see. It is therefore easy for a practitioner to change
the post-processing parameters to select a desirable balance between
performance and non-disparity. \gp{Copied directly from the rebuttal. Does it belong here?}

However, in some applications, the cost of randomizing predictions for the
stronger group ($G_2$) may not be a suitable option. In such cases, the only path towards
non-discrimination is improving the classifier for the weaker group ($G_1$)
either by collecting more data, using additional features, or denoising existing
data. When this extra work is not possible, achieving non-discrimination will
require a sacrifice in performance, as is the case with other frameworks
of fairness
\cite{berk2017fairness,corbett2017algorithmic,hardt2016equality,zliobaite2015relation}.
\gp{Trying to hammer the point home here, but a bit overkill maybe?}
In these settings, \autoref{alg:method} is
the best solution to achieve non-discrimination with calibration.

%such that $\epsilon(h_1') \le
%\epsilon(h_1)$
%and $\fp{h_1'} < \fp{h_1}$ or $\fn{h_1'} < \fn{h_1}$.
